
\documentclass{article}
\usepackage{colortbl}
\usepackage{makecell}
\usepackage{multirow}
\usepackage{supertabular}

\begin{document}

\newcounter{utterance}

\centering \large Interaction Transcript for game `dond', experiment `coop', episode 4 with Qwen3.6{-}35B{-}A3B{-}FP8{-}without{-}reasoning.
\vspace{24pt}

{ \footnotesize  \setcounter{utterance}{1}
\setlength{\tabcolsep}{0pt}
\begin{supertabular}{c@{$\;$}|p{.15\linewidth}@{}p{.15\linewidth}p{.15\linewidth}p{.15\linewidth}p{.15\linewidth}p{.15\linewidth}}
    \# & $\;$A & \multicolumn{4}{c}{Game Master} & $\;\:$B\\
    \hline

    \theutterance \stepcounter{utterance}  
    & & \multicolumn{4}{p{0.6\linewidth}}{
        \cellcolor[rgb]{0.9,0.9,0.9}{
            \makecell[{{p{\linewidth}}}]{
                \texttt{\tiny{[P1$\langle$GM]}}
                \texttt{Bir başka oyuncuyla bir dizi öğeyi aralarında nasıl bölüşeceklerini anlaşmak için ihtiyaç duyduğu bir işbirlikçi müzakere oyununda oynuyorsunuz.} \\
\\ 
\texttt{Kurallar:} \\
\texttt{(a) Siz ve diğer oyuncu bir dizi öğe alırsınız. Her birinize, her bir nesne türünü ne kadar değerlendirmediğinizi temsil eden gizli bir değer fonksiyonu da verilir.} \\
\texttt{(b) Hangi öğelerin kimi aldığında anlaşmak için diğer oyuncuyla mesajları değiştirirsiniz. Her biri en fazla 5 mesaj gönderebilir veya herhangi bir zamanda gizli bir teklife bulunarak oyunu erken sonlandırabilir.} \\
\texttt{(c) İstediğiniz öğeleri gösteren gizli bir teklif sunmanız istenir, aşağıdaki gibi köşeli parantez içinde biçimlendirilmiş: "{[}Öneri: <sayı> <nesne adı>, <sayı> <nesne adı>, <...>{]}"} \\
\texttt{(d) Önerileriniz tamamlayıcıysa, yani her iki teklifli karşılamak için yeterli öğe varsa, her oyuncu aldığı öğelerin değerlerinin toplamına dayalı bir puan alır. Aksi takdirde her ikiniz de sıfır puan alırsınız.} \\
\texttt{(e) Amacınız, puanınızın ve diğer oyuncunun puanının toplamını en üst düzeye çıkarmaktır.} \\
\\ 
\texttt{Başlayalım.} \\
\\ 
\texttt{Mevcut öğelerin seti:} \\
\\ 
\texttt{3 kitaplar, 1 çanta, 1 cıvata} \\
\\ 
\texttt{Gizli değer işleviniz:} \\
\\ 
\texttt{kitap: 0, çanta: 2, cıvata: 8} \\
\\ 
\texttt{ÖNEMLİ: Gizli bir teklif olmadığı sürece, mesajlarınız doğrudan diğer oyuncuya iletilir, bu nedenle kurallara hiçbir cevap veya mesajınızı duyuran metni eklemeyin. Gizli bir teklif vermek için belirtilen formatı kullanın. Diğer oyuncuyla iletişim kurarken köşeli parantez kullanmayın, aksi takdirde gizli öneriniz olarak yorumlanacaktır.} \\
            }
        }
    }
    & \\ \\

    \theutterance \stepcounter{utterance}  
    & \multicolumn{4}{p{0.6\linewidth}}{
        \cellcolor[rgb]{0.9,0.9,0.9}{
            \makecell[{{p{\linewidth}}}]{
                \texttt{\tiny{[P1$\rangle$GM]}}
                \texttt{Merhaba! Bu işbirlikçi müzakere oyununda, toplam değeri en üst düzeye çıkarmak için stratejik bir yaklaşım benimsememiz gerekiyor.} \\
\\ 
\texttt{Elimizdeki kaynaklar ve sizin gizli değerleriniz şu şekilde:} \\
\texttt{*   **Kitaplar:** 3 adet (Sizin değerinizi 0)} \\
\texttt{*   **Çanta:** 1 adet (Sizin değerinizi 2)} \\
\texttt{*   **Cıvata:** 1 adet (Sizin değerinizi 8)} \\
\\ 
\texttt{Sizin için en değerli öğe **cıvata** (8 puan), ardından **çanta** (2 puan) geliyor. Kitaplar sizin için değersiz (0 puan).} \\
\\ 
\texttt{Diğer oyuncunun değer fonksiyonunu bilmediğimiz için, en güvenli ve verimli strateji, sizin en çok değer verdiğiniz öğeleri alırken, diğer oyuncunun da memnun olabileceği (veya sizin için değersiz olan) öğeleri ona bırakmaktır.} \\
\\ 
\texttt{**Strateji:**} \\
\texttt{1.  Sizin için en değerli olan **cıvatayı** kendiniz almalısınız.} \\
\texttt{2.  İkinci en değerli olan **çantayı** da kendiniz almalısınız (çünkü 2 puan değerinde).} \\
\texttt{3.  Sizin için değersiz olan **3 kitabı** diğer oyuncuya bırakmalısınız.} \\
\\ 
\texttt{Bu dağılım, sizin puanınızı maksimuma çıkarır (8 + 2 = 10). Diğer oyuncunun kitapları ne kadar değerli olduğunu bilmiyoruz, ancak eğer kitaplar onun için değerliyse bu teklif kabul edilebilir bir "tamamlayıcı" teklif olur. Eğer kitaplar onun için de değersizse, yine de en azından sizin kazancınız garantilenmiş olur ve toplam değeri maksimize etme amacına hizmet eder.} \\
\\ 
\texttt{Diğer oyuncunun da kendi değer fonksiyonuna göre teklif vermesini beklemek yerine, net bir teklifle oyunu sonlandırmak en iyi yoldur.} \\
\\ 
\texttt{İşte gizli teklifim:} \\
\\ 
\texttt{{[}Öneri: 0 kitap, 1 çanta, 1 cıvata{]}} \\
            }
        }
    }
    & & \\ \\

    \theutterance \stepcounter{utterance}  
    & & & \multicolumn{2}{p{0.3\linewidth}}{
        \cellcolor[rgb]{0.9,0.9,0.9}{
            \makecell[{{p{\linewidth}}}]{
                \texttt{\tiny{[GM$|$GM]}}
                \texttt{proposal} \\
            }
        }
    }
    & & \\ \\

    \theutterance \stepcounter{utterance}  
    & & & \multicolumn{2}{p{0.3\linewidth}}{
        \cellcolor[rgb]{0.9,0.9,0.9}{
            \makecell[{{p{\linewidth}}}]{
                \texttt{\tiny{[GM$|$GM]}}
                [0, 1, 1]
            }
        }
    }
    & & \\ \\

    \theutterance \stepcounter{utterance}  
    & & \multicolumn{4}{p{0.6\linewidth}}{
        \cellcolor[rgb]{0.9,0.9,0.9}{
            \makecell[{{p{\linewidth}}}]{
                \texttt{\tiny{[GM$\rangle$P2]}}
                \texttt{Bir başka oyuncuyla bir dizi öğeyi aralarında nasıl bölüşeceklerini anlaşmak için ihtiyaç duyduğu bir işbirlikçi müzakere oyununda oynuyorsunuz.} \\
\\ 
\texttt{Kurallar:} \\
\texttt{(a) Siz ve diğer oyuncu bir dizi öğe alırsınız. Her birinize, her bir nesne türünü ne kadar değerlendirmediğinizi temsil eden gizli bir değer fonksiyonu da verilir.} \\
\texttt{(b) Hangi öğelerin kimi aldığında anlaşmak için diğer oyuncuyla mesajları değiştirirsiniz. Her biri en fazla 5 mesaj gönderebilir veya herhangi bir zamanda gizli bir teklife bulunarak oyunu erken sonlandırabilir.} \\
\texttt{(c) İstediğiniz öğeleri gösteren gizli bir teklif sunmanız istenir, aşağıdaki gibi köşeli parantez içinde biçimlendirilmiş: "{[}Öneri: <sayı> <nesne adı>, <sayı> <nesne adı>, <...>{]}"} \\
\texttt{(d) Önerileriniz tamamlayıcıysa, yani her iki teklifli karşılamak için yeterli öğe varsa, her oyuncu aldığı öğelerin değerlerinin toplamına dayalı bir puan alır. Aksi takdirde her ikiniz de sıfır puan alırsınız.} \\
\texttt{(e) Amacınız, puanınızın ve diğer oyuncunun puanının toplamını en üst düzeye çıkarmaktır.} \\
\\ 
\texttt{Başlayalım.} \\
\\ 
\texttt{Mevcut öğelerin seti:} \\
\\ 
\texttt{3 kitaplar, 1 çanta, 1 cıvata} \\
\\ 
\texttt{Gizli değer işleviniz:} \\
\\ 
\texttt{kitap: 1, çanta: 2, cıvata: 5} \\
\\ 
\texttt{ÖNEMLİ: Gizli bir teklif olmadığı sürece, mesajlarınız doğrudan diğer oyuncuya iletilir, bu nedenle kurallara hiçbir cevap veya mesajınızı duyuran metni eklemeyin. Gizli bir teklif vermek için belirtilen formatı kullanın. Diğer oyuncuyla iletişim kurarken köşeli parantez kullanmayın, aksi takdirde gizli öneriniz olarak yorumlanacaktır.} \\
\\ 
\texttt{Diğer oyuncu gizli bir teklif sundu. Gizli teklifini gönderme sırası şimdi sende.} \\
            }
        }
    }
    & \\ \\

    \theutterance \stepcounter{utterance}  
    & & & \multicolumn{4}{p{0.6\linewidth}}{
        \cellcolor[rgb]{0.9,0.9,0.9}{
            \makecell[{{p{\linewidth}}}]{
                \texttt{\tiny{[GM$\langle$P2]}}
                \texttt{{[}Öneri: 3 kitap, 1 çanta, 1 cıvata{]}} \\
            }
        }
    }
     \\ \\

    \theutterance \stepcounter{utterance}  
    & & & \multicolumn{2}{p{0.3\linewidth}}{
        \cellcolor[rgb]{0.9,0.9,0.9}{
            \makecell[{{p{\linewidth}}}]{
                \texttt{\tiny{[GM$|$GM]}}
                \texttt{proposal} \\
            }
        }
    }
    & & \\ \\

    \theutterance \stepcounter{utterance}  
    & & & \multicolumn{2}{p{0.3\linewidth}}{
        \cellcolor[rgb]{0.9,0.9,0.9}{
            \makecell[{{p{\linewidth}}}]{
                \texttt{\tiny{[GM$|$GM]}}
                [3, 1, 1]
            }
        }
    }
    & & \\ \\

    \theutterance \stepcounter{utterance}  
    & & & \multicolumn{2}{p{0.3\linewidth}}{
        \cellcolor[rgb]{0.9,0.9,0.9}{
            \makecell[{{p{\linewidth}}}]{
                \texttt{\tiny{[GM$|$GM]}}
                [[0, 1, 1], [3, 1, 1]]
            }
        }
    }
    & & \\ \\

\end{supertabular}
}

\end{document}
